%% ma: L11-B14-0027
%% lop: 11
%% chuong: 4
%% bai: 14
%% dang: 11.14.4
%% loai: TLN
%% muc: TH
%% dap_an: "0,33"
%% nguon: "tai-lieu-hoc-tap-mon-toan-lop-11-theo-chuong-trinh-giao-duc-pho-thong-2018.pdf (Lớp Toán Thầy Duy, Phạm Lê Duy, THPT Chu Văn An), Chương 4 Bài 5, Phần 3, Câu 2 (tr.535)"
%% kien_thuc: "Lớp 11 Bài 14 (phép chiếu song song)"
%% trang_thai: cach-ly
%% ly_do: "Đề gốc không nêu điểm E (chỉ xuất hiện trong câu hỏi EN/ED); đáp số 0,33 chỉ đúng nếu E là trung điểm AB. Cần giáo viên bổ sung 'Gọi E là trung điểm của AB' trước khi nhập."
%% ghi_chu: "Đã kiểm bằng tọa độ: N = (0,4; 0,1; 1/3) nằm trên ED với EN/ED = 1/3."
\begin{ex}
Cho tứ diện $ABCD$, $M$ là trọng tâm của tam giác $ABC$. Gọi $N$ là hình chiếu song song của điểm $M$ theo phương $CD$ lên mặt phẳng $(ABD)$. Khi đó $\dfrac{EN}{ED}$ bằng bao nhiêu? (làm tròn đến hàng phần trăm)
\shortans{0,33}
\loigiai{\sloppy Gọi $E$ là trung điểm của $AB$. Theo phương $CD$ lên $(ABD)$: $C\mapsto D$, $E\mapsto E$, nên $M\in CE$ với $\dfrac{EM}{EC}=\dfrac13$ có hình chiếu $N\in ED$ với $\dfrac{EN}{ED}=\dfrac13\approx0{,}33$.}
\end{ex}
